%%%%%%%%%%%%%%%%%%%%%%%%%%%%%%%%%%%%%%%%%%%%%%%%%%%%%%%%%%%%%%%%%%%%%%
\section{Travail proposé durant l'été}
%%%%%%%%%%%%%%%%%%%%%%%%%%%%%%%%%%%%%%%%%%%%%%%%%%%%%%%%%%%%%%%%%%%%%%

Le projet ALORA étant conçu comme un projet de recherche à long terme,
il apparaît nécessaire de commencer par la construction de fondations
solides avant d'aborder le développement des modèles d'intelligence
artificielle.

L'objectif de cette première phase n'est donc pas de développer un
modèle de reconnaissance vocale ou de traduction, mais de construire
une infrastructure scientifique permettant de collecter, organiser et
préparer les données nécessaires aux travaux futurs.

%%%%%%%%%%%%%%%%%%%%%%%%%%%%%%%%%%%%%%%%%%%%%%%%%%%%%%%%%%%%%%%%%%%%%%

\subsection{Objectifs de cette première phase}

Les objectifs principaux sont les suivants :

\begin{enumerate}

\item Définir une méthodologie rigoureuse de collecte de données
orales.

\item Concevoir une plateforme simple permettant la participation de
locuteurs natifs.

\item Définir la structure du corpus ainsi que les métadonnées à
conserver.

\item Concevoir un pipeline de validation et de contrôle qualité des
enregistrements.

\item Produire un premier corpus exploitable pour les expérimentations
menées durant le semestre S5.

\end{enumerate}

%%%%%%%%%%%%%%%%%%%%%%%%%%%%%%%%%%%%%%%%%%%%%%%%%%%%%%%%%%%%%%%%%%%%%%

\subsection{Plateforme de collecte}

Une plateforme web légère sera développée afin de permettre à des
locuteurs natifs d'enregistrer des mots ou des phrases en soussou.

Le principe envisagé est le suivant :

\begin{itemize}

\item présentation d'un mot ou d'une phrase en français ;

\item enregistrement de la traduction orale par le participant ;

\item possibilité de réenregistrer avant validation ;

\item enregistrement automatique des métadonnées.

\end{itemize}

Cette approche permettra de constituer progressivement un corpus
contenant plusieurs réalisations d'une même phrase par différents
locuteurs.

%%%%%%%%%%%%%%%%%%%%%%%%%%%%%%%%%%%%%%%%%%%%%%%%%%%%%%%%%%%%%%%%%%%%%%

\subsection{Métadonnées envisagées}

Chaque enregistrement devra être accompagné d'informations permettant
de faciliter les analyses futures.

À titre indicatif, les métadonnées suivantes sont envisagées :

\begin{itemize}

\item identifiant du locuteur (anonyme) ;

\item âge (optionnel) ;

\item sexe (optionnel) ;

\item région d'origine ;

\item langue maternelle ;

\item langues parlées ;

\item date d'enregistrement ;

\item qualité estimée de l'audio ;

\item texte français correspondant ;

\item durée de l'enregistrement.

\end{itemize}

Cette liste reste volontairement ouverte et pourra évoluer selon les
recommandations du laboratoire.

%%%%%%%%%%%%%%%%%%%%%%%%%%%%%%%%%%%%%%%%%%%%%%%%%%%%%%%%%%%%%%%%%%%%%%

\subsection{Organisation des données}

Une attention particulière sera portée à l'organisation du corpus afin
de faciliter son exploitation.

Les travaux prévus comprennent notamment :

\begin{itemize}

\item définition de l'arborescence des données ;

\item choix des formats audio ;

\item convention de nommage des fichiers ;

\item organisation des métadonnées ;

\item préparation d'un format facilement exploitable par les outils de
Machine Learning.

\end{itemize}

%%%%%%%%%%%%%%%%%%%%%%%%%%%%%%%%%%%%%%%%%%%%%%%%%%%%%%%%%%%%%%%%%%%%%%

\subsection{Contrôle qualité}

Avant toute utilisation scientifique, les données devront être
contrôlées afin de garantir leur qualité.

Plusieurs mécanismes sont envisagés :

\begin{itemize}

\item détection des enregistrements silencieux ;

\item contrôle de la durée ;

\item détection de fichiers corrompus ;

\item suppression des doublons ;

\item vérification de la cohérence entre les métadonnées et
l'enregistrement.

\end{itemize}

%%%%%%%%%%%%%%%%%%%%%%%%%%%%%%%%%%%%%%%%%%%%%%%%%%%%%%%%%%%%%%%%%%%%%%

\subsection{Résultats attendus}

À l'issue de cette première phase, les livrables attendus sont les
suivants :

\begin{enumerate}

\item une plateforme fonctionnelle de collecte de données ;

\item une première version du corpus oral ;

\item une documentation décrivant la structure des données ;

\item un pipeline de validation et de préparation du corpus ;

\item un ensemble de recommandations méthodologiques pour les étapes
suivantes du projet.

\end{enumerate}

%%%%%%%%%%%%%%%%%%%%%%%%%%%%%%%%%%%%%%%%%%%%%%%%%%%%%%%%%%%%%%%%%%%%%%

\begin{discussion}

Cette première phase représente principalement un travail de
construction méthodologique.

Je souhaiterais profiter de cette période pour bénéficier de votre
expérience afin de valider les choix scientifiques concernant la
collecte des données, l'organisation du corpus, les métadonnées à
conserver ainsi que les outils les plus adaptés avant le démarrage des
travaux de modélisation.

\end{discussion}